\section{Mathematical Foundations of the Translation Mechanism}
\label{sec:mathfound}

The tree-to-graph-to-tree pipeline of Section~\ref{sec:ontology-bridge} is not an ad hoc
engineering convenience; it is an instance of a construction that recurs, under different
names, across category theory, applied topology, and computational linguistics. Making the
construction explicit serves two purposes: it shows \emph{why} the round trip resolves the
N:M mappings that defeat direct tree-to-tree approaches, and it makes the frontier of the
method legible --- it tells us exactly where deterministic guarantees hold and where
irreducible reconciliation loss begins.

\subsection{The Babel Problem, Stated Formally}
Let a localized statutory chart of accounts be a rooted tree $T_s$ and a target statutory
chart be a rooted tree $T_t$, each carrying the single-parent constraint of
Section~\ref{sec:ontology-bridge}. The ``Accounting Babel'' is the observation that there is
in general \emph{no faithful structure-preserving map} $T_s \rightarrow T_t$: an account in
$T_s$ may correspond to several accounts in $T_t$ and vice versa (the $1{:}N$ and $N{:}M$
cardinalities), and the multi-dimensional roles an account plays (statutory, fiscal,
functional, consolidated) cannot be expressed inside a single-parent hierarchy at all. The
direct approach fails because the two trees do not share a common coordinate system.

\emph{A concrete instance.} Consider one economic event --- a domestic credit
sale of goods with value-added tax. In Mexico's SAT chart it touches
\texttt{105} (Clientes), \texttt{401} (Ingresos), and \texttt{208}
(IVA trasladado); in France's PCG, \texttt{411} (Clients), \texttt{707}
(Ventes de marchandises), and \texttt{44571} (TVA collect\'ee); in Vietnam's
VAS, \texttt{131} (trade receivables), \texttt{511} (revenue), and
\texttt{3331} (VAT payable). No single-parent edge relates ``SAT 208'' to
``PCG 44571'' to ``VAS 3331'': they are the \emph{same} concept
(output VAT payable) under three incompatible coordinate systems, and the
receivable line additionally carries fiscal and functional roles the tree
cannot hold at once. This is the $N{:}M$ obstruction the formalism names ---
and the reason the pivot below routes every chart through one shared node set
rather than attempting $T_s \rightarrow T_t$ directly.

\subsection{Lift, Pivot, Project: an Adjoint Pair}
Kontablo resolves this exactly as machine translation resolves the analogous problem between
natural languages: rather than mapping $T_s$ directly to $T_t$, it analyzes $T_s$ up into a
jurisdiction-independent intermediate representation and then generates $T_t$ down from it.
In machine translation this intermediate is the \textit{interlingua} at the apex of the
Vauquois triangle (Vauquois, 1968); in Kontablo it is the Level~3 universal node set
$V_{\text{universal}}$, which serves as the absolute pivot.

This pivot construction has a natural categorical reading. Following Spivak (2012), a
database schema is a small category and an instance is a set-valued functor on it; a schema
mapping $F: \mathcal{S} \rightarrow \mathcal{T}$ then induces three canonical
\emph{data-migration functors}. Kontablo's import and export streams are structurally
analogous to two of them:
\[
\Sigma_F : \mathcal{S}\text{-Inst} \rightarrow \mathcal{T}\text{-Inst}
\qquad\text{(import / lift: tree}\rightarrow\text{graph)}
\]
\[
\Delta_F : \mathcal{T}\text{-Inst} \rightarrow \mathcal{S}\text{-Inst}
\qquad\text{(export / project: graph}\rightarrow\text{tree)}
\]
where $\Sigma_F$ is left adjoint to $\Delta_F$. The import stream enriches each leaf of the
local tree with a mapping edge to a universal node and lifts it into the multi-dimensional
graph $G$; the export stream linearizes $G$ back to a single-parent tree. We state the
correspondence as an analogy, in the same register as the Cover's-theorem passage below,
because the full construction has not been carried out: we have not defined the schema
categories, constructed $F$, or verified the adjunction, and the implemented export step
\emph{selects} a primary reporting dimension (Section~\ref{sec:ontology-bridge}) --- a
choice a canonical adjoint would have to eliminate or parameterize. Working the
correspondence out for a minimal chart pair, or amending it where it fails, is labeled
future work; what the analogy buys today is a design discipline (imports should be free
constructions, exports should be structure-forgetting) rather than a theorem.

\subsection{Why Lifting the Dimension Resolves the Mapping}
The intuition that a problem unsolvable in a low-dimensional structure becomes solvable once
lifted into a richer one is formalized, in a different setting, by Cover's theorem
(Cover, 1965): a classification that is not separable in a low-dimensional space becomes
separable, with high probability, once mapped nonlinearly into a higher-dimensional space.
We invoke this only as an analogy of \emph{principle} --- Cover's setting is geometric
separability, not structure-preserving translation --- but the shared moral is exact: the
$N{:}M$ entanglement that cannot be resolved inside a single-parent tree is resolved once the
account is re-expressed as an edge into a multi-dimensional graph, where membership in
several analytical dimensions is no longer in conflict.

\subsection{The Cost of the Return Trip: a Reconciliation Obstruction}
\label{subsec:reconciliation_obstruction}
The return trip is lossy by construction. Linearization selects a primary reporting
dimension and discards the secondary analytical edges
(Section~\ref{sec:ontology-bridge}); the discarded content is precisely the failure of the
graph $G$ to be a tree. The natural question --- \emph{when can the local ledgers of many
jurisdictions be glued into one globally consistent account, and what obstructs it?} --- is
the question that sheaf cohomology was built to answer.

Assign to the mapping graph a cellular sheaf whose stalks carry the local ledger data and
whose restriction maps encode the consistency conditions on shared (overlapping) accounts
(Robinson, 2017; Curry, 2014; Hansen \& Ghrist, 2019). The domains here are explicit: each
stalk is a finite-dimensional real vector space $\mathbb{R}^{d}$ (the monetary amounts a node
carries across its $d$ analytical dimensions, valued in the Pacioli group of
Section~\ref{subsec:pacioli}), the restriction maps are $\mathbb{R}$-linear, and the
cohomology spaces are therefore $\mathbb{R}$-vector spaces. Then the degree-zero cohomology
$H^{0}$ is the space of \emph{globally consistent reconciliations} --- the data that patches
together across all jurisdictions --- while a non-trivial degree-one class, $H^{1} \neq 0$,
is the exact obstruction to such gluing: it measures, and localizes, the inconsistency that
no choice of mapping can remove.

This gives two of Kontablo's operational notions an intended formal referent. A node's
\emph{deterministic boundary} (ADR~009) is designed to play the role of a sheaf restriction
map: it constrains which local sections are admissible. The \emph{Inconsistency Flag} of
the co-responsibility architecture is designed to play the role of a witness to a local
section that fails to extend --- a non-zero cohomology class. The implemented flag is, today,
a set of deterministic rule checks (Section~\ref{subsec:expanded_validation}), not an
instantiated sheaf; the current state of formal validation is a committed toy experiment
(\texttt{research/experiments/kirchhoff\_hodge\_toy/}) that constructs a cellular sheaf over
a small intercompany scenario, computes $\delta^{0}$ and the $H^{0}$/$H^{1}$ dimensions
numerically, and resolves the imbalance by Moore--Penrose reconciliation --- a mechanics
validation, not a proof of the general correspondence. Determinism and reconciliation are
nonetheless designed to speak one language: the deterministic scaffolding defines the
restriction maps, and cohomology is the intended formal account of where global consistency
fails and a human must dispose.

\subsection{The Accounting-Native Anchor: the Pacioli Group}
\label{subsec:pacioli}
None of this is foreign to accounting's own mathematics. Ellerman (2014) shows that
double-entry bookkeeping is, exactly, the \emph{group of differences} construction (the
``Pacioli group'') on unsigned T-account pairs, and --- crucially for Kontablo --- that this
construction \emph{generalizes naturally to multi-dimensional vector accounting}, in which
scalar money amounts are replaced by vectors recording distinct types of property or
analytical dimension. This fixes the domain of monetary value precisely. The scalar Pacioli
group is the group of differences of ordered pairs of non-negative amounts,
$\mathcal{G} = (\mathbb{R}_{\geq 0} \times \mathbb{R}_{\geq 0})/\!\sim\, \cong (\mathbb{R}, +)$:
a signed real, where the debit/credit nature supplies the sign rather than a separate
$\mathbb{R}_{\geq 0}$ magnitude. The multi-dimensional generalization is $\mathcal{G}^{d}
\cong \mathbb{R}^{d}$ over the $d$ analytical dimensions --- the domain the sheaf stalks of
Section~\ref{subsec:reconciliation_obstruction} are valued in. (Implementations realize
$\mathcal{G}$ as fixed-point rationals $\subset \mathbb{Q}$ for exact cent arithmetic; the
algebra is stated over $\mathbb{R}$.) Kontablo's multi-dimensional property graph is the data-structure
realization of precisely that generalization. The contribution is therefore not to invent a
new algebra of accounting, but to operationalize, as a deterministic protocol over a
property graph, a multi-dimensional generalization that has been latent in the mathematics of
double-entry bookkeeping since its first formal treatment.
